\section{The Restoration of Aristocracy}

Aristocracy means rule by the best. Hence, anyone interested in the reconstitution of an elite must first understand what is the best way to live. Here we can start with some ideas from \textbf{Rene Guenon} to see how they clarify and expand on \textbf{Aristotle}'s ideas on living the good life. For Guenon there are two options:

\begin{enumerate}


\item Actualize all the possibilities of the human state

\item Actualize all the possibilities of transcendent states
\end{enumerate}


\textbf{Aristotle} was limited to the first of these, no mean accomplishment, so we shall start with that and conclude with the second option.

\paragraph{Possibilities of the Good Life}


For Aristotle, the pursuit of the good life was not an option for all. So much depends on good fortune, such as the possession of good health, intelligence, financial resources, a good family background, and so on. There must be an excess beyond the mere need to survive. Guenon would describe it in terms of one's possibilities arising from the caste the person would belong to. Christians instead ascribe such fortune to diving providence, predestination, or predilection.

For the purposes of this discussion we will leave aside the question about the possibilities of those who have not been blessed with the requisite resources Aristotle describes. The pagan in any case was indifferent to them, since all is in the hands of the fates or the whims of the gods. In Hinduism, for example, the Brahmans are unconcerned with the plight of the dalits. Even for Buddhism, which introduced the idea of compassion, the goal is not to relieve the external causes of suffering but rather to teach the elimination of desire as the root cause of suffering.

A curious superstition has arisen in our time, going by the name, the ``Law of Attraction''. This is a form of primitive theurgy that wishes to compel the universe to deliver the goods through thought or visualizations. However, this is defective in that it seeks only to receive the material conditions for the good life without undergoing the moral conversion to actually achieve it. The self is creating its image in consciousness all day long through its multitudinous thoughts and imaginations; spending a mere five or ten minutes trying to consciously concentrate one's thoughts and visualizations in a positive direction is insufficient to counteract the random perturbations arising in consciousness the rest of the day.

\paragraph{False Forms of the Good Life}


The first false form of the good life is to live solely for pleasures. There is little point to discuss this here. The other form is to live to achieve esteem in the community. This is a powerful motivating factor for those who are able to achieve it and may sometimes have a public good. In my visits to a hospital this week, I noticed all the annexes named for prominent or wealthy people in the local community. Their need for public esteem paid for special centers to treat cancer and other diseases.

In a disordered society, however, those held in esteem may be much less worthy such as the attention given to various entertainers of dubious moral standing. Now we even have those who are ``famous for being famous'', which fame they achieved comes from so-called ``reality'' shows, rather than for any worthwhile accomplishment.

Esteem can also be received on a lower scale, the so-called ``big fish in a small pond''. Thus, there are a myriad number of small groups promoting various political and religious causes. Their leaders publish books and are invited to speak at conferences. As we recently pointed out, there will always to supporters for some cause. However, the lower cannot judge the higher, so the quantity of supporters is a poor guide to determine the actual worth of those leaders. A solution will be proposed in the next section.

\paragraph{Promoting the Good Life}


The good life consists in actualizing all one's possibilities. Since the distinctive feature of the human being is the rational or intellectual soul, this can only mean to live one's life on that basis. Such a man will dominate the irrational elements in his soul and live by his intelligence. The intelligence will be guided by the Logos, or the rational order of things.

The true aristocrats, therefore, are those willing and able to follow that path. If such men can come together in a true friendship, the possibility of an elite can arise. There are two necessary elements, the proper education and a code of honour.

The aristocrat should have a grounding in liberal arts, specifically the trivium and the quadrivium. These teach the basics of logic, mathematics, language, and cosmology which are the foundation for advanced studies. As the potential rulers, they don't necessarily need to go deeply into philosophy and theology. The practical arts such as medicine, law, engineering are usually more appropriate for the producing class. The point is to be able to think rationally so they can follow an argument and be immune to false teachings. Through rhetoric, they learn the techniques of persuasion. Although the rulers must be willing to use force when necessary, persuasion is the better approach. Furthermore, with this training they will be able to recognize the rhetorical devices of false teachers and will not be swayed by them.

It would be premature to formulate a code of honour by comparing various ancient and medieval codes. The important point for now is how to enforce it. Since the aristocrat is the ``best'', they cannot be judged by those lower. Hence, they must be self-monitoring. This is where the moral will comes in because of the tendency to protect or promote one's friends, or even one's own fortune. Nevertheless, this elite would have to be willing to exclude those who violate that code.

In the novel, \emph{The Four Feathers} by \textbf{A E W Mason}, a man disgraces himself by leaving the army at a time of war. Three of his friends, as well as his fiancée, each give him a feather to point this out. This is the model that needs to be followed by groups with a political aspiration, as least those who do not count on the popularity of the crowd. They need to establish a code of honour and the will to enforce it. Those involved in an immoral lifestyle, the vulgar in rhetoric, the irrational, etc., must be excluded from leadership positions by group consensus. This we see nowhere, so their efforts will come to naught. Those of an aristocratic mind would withhold all support from such movements.

\paragraph{The Sage}


The philosopher is the lover of wisdom and is in search of it. For some, that search is a way of life. Hence, they engage in debates, disputes, arguments, dialogues, conferences, etc., that simply prolong the search without moving closer to the goal.

The Sage, on the other hand, is the one who has achieved wisdom. Beyond the aristocrat, these are the initiates, the saints, and so on, who have actualized possibilities beyond the human state. Their training is done in private, since it cannot be encapsulated fully in the written word.

Aristotle had no awareness of this level, but they must serve as the spiritual and intellectual guides to the aristocrats. Through the rites, moral codes, and exoteric religion they can connect with the folk, providing a common mind for the community to give it cohesion and continuity. Only they will have the authority to anoint the leaders and give them the legitimacy to lead effectively.


\flrightit{Posted on 2013-10-23 by Cologero}

\begin{center}* * *\end{center}

\begin{footnotesize}
\begin{sffamily}
\texttt{David on 2013-10-23 at 10:10 said:}

Very interesting. It goes along well with Boethius thoughts on the quadrivium needed for the aristocratic; and as Boethius was named the last real roman aristocrat, it suits well the discussion.

\hfill

\texttt{JA on 2013-10-23 at 10:31 said:}

I guess I'm not an aristocrat after all, I see now that I lack the qualifications -- I have some in that I come from a good heritage, had an excellent education, and am far above average intelligence but I due to events not in my power to control I lack financial security. Should I try to follow a path of a lower caste then ?

\hfill

\texttt{Michael on 2013-10-23 at 11:12 said:}

Would the law of attraction work with a moral person who has mastered the ``random perturbations arising in consciousness the rest of the day.'' If so, it would be a tool for achieving the good life described above.

It also seems that it is prudent for those who seek to transform society to achieve positions of leadership in the current society. For example, within academia, the Church, the military, and in business.

\hfill

\texttt{Caleb Cooper on 2013-10-23 at 11:15 said:}

Cologero, what do you think of the German system where the leading Princes/Cardinals elect the Emperor/Pope? Is it better or worse than hereditary monarchy?

At first glance it strikes me as better than hereditary monarchy, which suffers from regression to the mean. Having the highest peers choose to elevate one of their own seems to offers more safe guards, while still being unmistakeably and unapologitcally elitist. (I suppose you could have both if the Emperor was given special dispensation to have a `harem', which in the modern world of paternity testing wouldn't need to be a real cloistered harem (and no, I'm not projecting because I have no delusions about being an Emperor; I'm a simple yeoman, not an aristocrat with the presence to rule a group. Humans show a strong preference for dynasties and supermen though, so letting the Emperor have a huge number of children and electing the next Emperor from those who have at least 1/8th of his blood would be a large enough pool to find someone of capable quality))

On the topic of yeoman, do they have a place making decisions at the smaller local level? I view the practical quality of the aristocracy as possessing the confidence to effectively lead other men. Yeomen may not have as much leadership ability, but unlike the Shudra, they own themselves and their own labor. Maybe I'm trying to make too much of a sub distinction.

How large do you imagine the aristocracy being? Is there a place for a lower tier that would allow some of the yeomen and middle class to participate in local politics like electing the town administrator (mayor)? I suppose I really should study more on how the medieval world was run.

\hfill

\texttt{Cologero on 2013-10-23 at 12:08 said:}

JA, I don't believe I mentioned ``financial security'' anywhere; I'll retract it if you can find it. Obviously, you have sufficient leisure time and resources to engage in discussions such as this.

\hfill

\texttt{Cologero on 2013-10-23 at 12:13 said:}

Exactly, Michael, as a man thinketh in his heart, so he is.

And ditto for point 2, and a very Gramscian notion. It demonstrate initiative and leadership. Recently, there was a conference of some ``Traditional Britain'' group that I mentioned last week. They decided that their plan of action was to wait for some saviour to come and lead them out of the morass they believe they are in.

\hfill

\texttt{Cologero on 2013-10-23 at 12:22 said:}

Actually, that system was in place in the ancient Roman Kingdom era. Heredity may provide the proper conditions, such as access to education and preparation for leadership. Yet, it is hardly a sufficient condition. Ultimately, the Emperor must be anointed to be legitimate, so the spiritual authorities need the power of the veto. In the Roman kingdom, there was also the requirement for a divine omen. (Hence the idea of a lightning strike as a sign from God).

Of course the yeomen have the right to their own decisions. That was the purpose of the guild system. Ideally, and even Evola describes it this way, the principle of subsidiarity limits decisions to the lowest level necessary.

If you read older texts, most people simply did not even consider the possibility of participating in politics, no more than they would run their own hospitals ... it was not their expertise. However, medicine is more objective and the bad doctors are identifiable and can be weeded out. Not so for political leaders. Unless there is that code of honour among their peers, that is the weak point.

\hfill

\texttt{JA on 2013-10-23 at 20:36 said:}

I interpreted ``There must be an excess beyond the mere need to survive.'' as meaning that one must be wealthy enough to not have to work for a living to be of the elite, as was customary among the aristocracy of traditional Europe.

I've benefited alot from the equalising affects of modernity because through the Internet, public libraries, and used book dealers, I've been able to acquire knowledge that in olden times would have cost a fortune in instruction. Would I have had this knowledge had I been born back in the middle ages ? maybe as my family were nobles but who can say.

\hfill

\texttt{JA on 2013-10-23 at 20:39 said:}

De Maistre's vision was a Europe of monarchies under a Pope who had the power to check abuses. I think his vision is Traditional and the best system.

\hfill

\texttt{Cologero on 2013-10-24 at 00:31 said:}

Yes, JA, that is the paradox of modernity. There are so many resources and opportunities, yet so few true aristocrats.

\hfill

\texttt{laughingtorch on 2013-10-24 at 09:43 said:}

Aristotle would not have been allowed to discuss the transcendent state/Mysteries/initiation in his `lecture notes'. This is no doubt why he only logically concludes the unmoved mover. He was no doubt an initiate, and a Sage to Alexander. He was however forbidden to divulge thru vulgar and popular means the Mysteries.

\hfill

\texttt{laughingtorch on 2013-10-24 at 09:44 said:}

(aside from that this is a really good article, Colgero -- i love this website!!)

\hfill

\texttt{Cologero on 2013-10-24 at 13:08 said:}

That explanation sounds plausible, laughing torch. After all, as Guenon points out, Aristotle's metaphysics is close to some Vedic schools, not likely solely by accident. See, for example, his essay Nama-Rupa included in Studies in Hinduism. Plato, also, claimed that not everything was written down.

\hfill

\texttt{Sebastianus Lusitanus on 2020-07-07 at 19:32 said:}

Dear Cologero,

I am a new reader of Gornahoor and I am very happy for having found the website. Thanks for keeping that up and running.

I would like to know your assessment on this quote from Meditations on the Tarot, p. 238, where Tomberg seems to go against the organized formation of an elite:

``At the same time there is no more serious error than that of wanting to ``organise'' a community or fraternity which would be called to play either the role of an instrument of spiritual selection or election, or even the role of a spiritual elite. Par one can neither assume the function of election nor consider oneself as elect.''

But isn't that what the XXth century pre-Vatican II right-wing movements tried to accomplish? (with study and action groups; Jean Ousset's book ``L'Action'' comes to mind).

Thank you.

\hfill

\texttt{Cologero on 2020-07-08 at 05:14 said:}

@Sebastianus. Tomberg is not talking about a political movement. Read the story of the Transfiguration. Right away, Peter wanted to erect a monument. In effect, that would mean to organise the mystical experience, and make it openly available to the public. The esoteric centre chooses you, you don't choose it.

\hfill

\texttt{Taivasvaeltaja on 2023-10-23 at 08:12 said:}

``... Christian Esotericism, as much as all the spiritual systems, preaches brotherhood... but not that it's first messenger Jesus Christ with his radical message of compassion and the spirit of individuated brotherhood of awakened men and women, would have taught that all human beings are equal... there are viewpoints that are more ethical and enlightened, more logical and reasonable and just, and then there are errors of judgment and even cruelty... therefore, Christianity, as I see it, is also aristocratic and not at all a democratic... but the spiritual adepts, the brothers of light, never rule upon their subjects with violence and oppression -- like the material aristocracies have done as failed human systems due to human fallibility in the sub-lunar realms -- but only with the gentle persuasion to make them understand true freedom and liberty from the fallen world... ''

-- Aarni Kouta, Lusiferin Kannel (with my own additions and thoughts)

\end{sffamily}
\end{footnotesize}

